%% ===========================================================================
%% Standalone fragment: §3.4 The Invariance-Blindness Theorem (IBT).
%% NOT wired into NOETHER_paper_arxiv.tex yet (no LaTeX toolchain in container;
%% B2 integration + compile audit per CLAUDE.md §8 to be done in a LaTeX env).
%% Uses the paper's existing environments: definition, lemma, theorem, remark,
%% corollary*. Cross-refs to be resolved on integration (marked % TODO-ref).
%% Source of truth: docs/tosem_maturity_2026-06-16/invariance_blindness_theorem_draft.md
%% ===========================================================================

\subsection{The Invariance-Blindness Theorem}
\label{subsec:ibt}

Theorem~\ref{thm:closure} guarantees that \texttt{CONSTRUCT-MP} drops no
\texttt{Translate}-reachable MR; it is, by construction, silent about what the
derived MRs can \emph{detect}. This subsection supplies that missing content: a
\emph{limiting} characterization of the faults an algebra-induced MR cannot
detect. Throughout we work in the linear operator-implementation fault class
(Definition~\ref{def:fault-class}); scope and assumptions are collected in
Remark~\ref{rem:ibt-scope}.

\begin{definition}[Linear program and fault class]
\label{def:fault-class}
A \emph{linear program} is $P_L:\mathbb{R}^N\to\mathbb{R}^N$, $P_L(x)=Lx$, indexed by an
operator $L\in\Theta:=\mathbb{R}^{N\times N}$ ($\dim\Theta=N^2<\infty$). The reference
(correct) program is $P_{L^{*}}$; an \emph{operator-implementation fault} is
$L=L^{*}+\Delta$ with $\Delta\in\Theta$. (For nonlinear solvers, $\Theta$ is the
operator's finite parameter vector and $E_s$ below its first-order linearization.)
\end{definition}

\begin{definition}[Structural defect and detection kernel]
\label{def:defect-kernel}
Let block $s\in\mathcal{D}(\mathcal{A}_P)$ carry a structure $T_s$ (for $s=G$, a group
action $A\in G\le GL(\mathbb{R}^N)$; for $s=T^{*}$, the inner-product adjoint). The
algebra-induced MR $\rho_{\iota,s}=\texttt{Translate}(\iota,s)$ asserts a structural
predicate $\Pi_s[P_L]$ over a witness set $W=S_s\times X_0$ ($S_s\subseteq T_s$,
$X_0\subseteq\mathbb{R}^N$). For $s=G$ the per-witness \emph{defect} is the commutator
applied to an input,
\[
E_G(L;A,x)\;=\;P_L(Ax)-A\,P_L(x)\;=\;[L,A]\,x,
\]
linear in $L$; for $s=T^{*}$, $E_{T^{*}}(L)=L-L^{\top}$. The \emph{detection kernel}
is the set of faults the MR fails to flag,
\[
\ker(\rho_{\iota,s})\;:=\;\{\,L\in\Theta : E_s(L;w)=0\ \ \forall w\in W\,\}.
\]
If $X_0$ spans $\mathbb{R}^N$ then $\ker(\rho_{\iota,s})=\bigcap_{A\in S_s}\mathcal{C}(A)$
with $\mathcal{C}(A)=\{L:[L,A]=0\}$ the commutant of $A$.
\end{definition}

\begin{definition}[Structure-preserving fault]
\label{def:compatible}
$L$ is \emph{$T_s$-compatible} iff it satisfies the full structural identity that
validates $\rho_{\iota,s}$: $E_s(L;w)=0$ for all witnesses $w\in W^{\mathrm{full}}_s$
(for $s=G$, $[L,A]=0$ for all $A\in T_s$). Write $\mathcal{C}(T_s)$ for the
$T_s$-compatible set (the commutant of the whole group, i.e.\ the equivariant
operators).
\end{definition}

\begin{definition}[Faithfulness]
\label{def:faithful}
A finite witness set $W^{\mathrm{test}}_s$ is \emph{faithful} iff
$\ker(\rho_{\iota,s})=\mathcal{C}(T_s)$; equivalently, since $E_s$ is linear in $L$,
$\operatorname{rank}J_{\mathrm{test}}=\operatorname{rank}J_{\mathrm{full}}$, where $J$ collects the
rows $L\mapsto E_s(L;w)$.
\end{definition}

\begin{lemma}[Reachability of faithfulness]
\label{lem:reachability}
If $\Theta$ is finite-dimensional and $E_s$ is linear in $L$, then a finite
faithful witness set $W^{\mathrm{test}}_s$ exists.
\end{lemma}

\begin{proof}
Each witness $w$ contributes a linear functional $a_w:L\mapsto E_s(L;w)\in\Theta^{*}$.
The span $V_{\mathrm{full}}=\operatorname{span}\{a_w:w\in W^{\mathrm{full}}_s\}\subseteq\Theta^{*}$ is a
subspace of the finite-dimensional dual, hence finite-dimensional; pick finitely
many witnesses $W^{\mathrm{test}}_s$ whose functionals span $V_{\mathrm{full}}$. As
$W^{\mathrm{test}}_s\subseteq W^{\mathrm{full}}_s$ the intercepts agree, so the two systems
have the same solution set, i.e.\ $\ker(\rho_{\iota,s})=\mathcal{C}(T_s)$.
\end{proof}

Lemma~\ref{lem:reachability} is the crux: tightness is \emph{not} tautological but
follows from finite fault-dimensionality, so a finite executable MR can pin the
exact blind set. We now characterize that set.

\begin{theorem}[Invariance-Blindness, $G$ and $T^{*}$ blocks]
\label{thm:ibt}
Let the fault class be linear (Definition~\ref{def:fault-class}), let $E_s$ be
linear in $L$ (holds for $s\in\{G,T^{*}\}$), and let the test witnesses be faithful
(Definition~\ref{def:faithful}). Then
\[
\ker(\rho_{\iota,s})\;=\;\mathcal{C}(T_s)\;=\;\{\,L:\ P_L\ \text{preserves}\ T_s\,\}.
\]
That is, a linear operator-implementation fault survives the structural MR
\emph{if and only if} it preserves the structure $T_s$ that the MR exploits.
\end{theorem}

\begin{proof}
($\supseteq$) If $L\in\mathcal{C}(T_s)$ then $E_s(L;w)=0$ on $W^{\mathrm{full}}_s\supseteq
W^{\mathrm{test}}_s$, so $\rho_{\iota,s}$ passes and $L\in\ker$.
($\subseteq$) If $L\in\ker$ then $E_s(L;w)=0$ on $W^{\mathrm{test}}_s$; by faithfulness
(Definition~\ref{def:faithful}) $\ker(\rho_{\iota,s})=\mathcal{C}(T_s)$, so
$L\in\mathcal{C}(T_s)$.
\end{proof}

\begin{remark}[Worked instance: advection--diffusion, translation block]
\label{rem:ibt-advdiff}
On a periodic grid the constant-coefficient operator $c\cdot\nabla+\alpha\nabla^2$
commutes with every grid translation for \emph{all} constants $(c,\alpha)$; hence
$E_G(L;\text{shift})\equiv 0$ on the whole constant-coefficient family, so
$\ker(\rho_{\iota,G})$ contains it, including the advection-speed fault $c\mapsto c'$.
The translation MR is therefore \emph{blind} to advection-speed errors, matching
the empirical $0/n$ of \S\ref{sec:empirical-evaluation}.% TODO-ref empirical sec
A spatially inhomogeneous coefficient $c(x)$ breaks translation equivariance, so
$E_G\neq 0$ and the fault is detected. A faithfulness check (one generator plus an
input basis is unisolvent) is reported in the supplement.% TODO-ref supplement S10
\end{remark}

%% corollary* is unnumbered (paper preamble: \newtheorem*{corollary*}). For
%% numbered, \ref-able corollaries, add \newtheorem{corollary}{Corollary} to the
%% preamble at integration and drop the stars.
\begin{corollary*}[Single-block incompleteness]
For any block $s$ with $\mathcal{C}(T_s)\supsetneq\{L^{*}\}$ (a nontrivial structure-
preserving fault exists), the single-block battery $\{\rho_{\iota,s}\}_\iota$ is
incomplete: its kernel contains that fault.
\end{corollary*}

\begin{corollary*}[Completeness requires a trivial joint kernel]
A family of oracles $\{O_j\}$ detects every nontrivial fault only if
$\bigcap_j\ker(O_j)=\{L^{*}\}$. Completeness therefore requires oracles whose
structural kernels intersect trivially, not additional MRs of the same block.
\end{corollary*}

\begin{corollary*}[Differential oracle is the complementary blind spot]
The neutral cross-implementation differential oracle $O_{\mathrm{diff}}$, which flags
a fault when two independent implementations disagree beyond a discretisation-
calibrated tolerance, has kernel $\ker(O_{\mathrm{diff}})=\{$common-mode faults
(identical effect on both implementations)$\}$. Since common-mode and
$T_s$-compatible are distinct conditions, $\ker(\rho_{\iota,s})\cap\ker(O_{\mathrm{diff}})$
is strictly smaller than either: the MR battery and the differential oracle are
complementary, and a fault escapes both only if it is simultaneously
structure-preserving and common-mode.
\end{corollary*}

\begin{remark}[Scope and assumptions]
\label{rem:ibt-scope}
(\textbf{R1}) Theorem~\ref{thm:ibt} is stated for the linear operator-implementation
fault class (Definition~\ref{def:fault-class}); it does not extend to arbitrary
nonlinear faults. (\textbf{R2}) Detection is exact-arithmetic; an executable MR with
tolerance $\tau$ has kernel $\ker_\tau\supseteq\ker$, so faults below $\tau$ are
false negatives. Theorem~\ref{thm:ibt} is the $\tau\to0$ limit; the finite-$\tau$
regime is the detectability floor of \S\ref{sec:empirical-evaluation}.% TODO-ref
(\textbf{R3}) Linearity of $E_s$ holds for $G$ and $T^{*}$ (commutator; symmetric
part). It fails for $O_\le$ (monotonicity is an inequality, a cone not a subspace),
for $\mathcal{T}^{*}_{\mathrm{rev}}$ (reversibility $R L R^{-1}-L^{-1}$ involves the inverse), and
for $\mathcal{L}^{*}$ (Richardson ratio of norms); for these blocks only the sufficient
direction ($\supseteq$) holds, or a linearized subclass. (\textbf{R4}) The sufficient
direction is by definition; the content is the $\subseteq$ direction together with
Lemma~\ref{lem:reachability} and the kernel characterization.
\end{remark}
